% Geometric nucleon binding and beta decay (HQIV) — working copy
% Self-contained mathematical narrative; machine-checked audit pointers in Appendix A only.

\documentclass[11pt,a4paper]{article}

\usepackage{amsmath,amssymb,amsthm,mathtools}
\usepackage{microtype}
\usepackage[hidelinks]{hyperref}
\usepackage{xurl}
\newcommand{\repo}{\url{https://github.com/disregardfiat/hqiv-lean}}
\usepackage[margin=1in]{geometry}

\theoremstyle{plain}
\newtheorem{theorem}{Theorem}
\newtheorem{lemma}[theorem]{Lemma}
\newtheorem{proposition}[theorem]{Proposition}
\newtheorem{corollary}[theorem]{Corollary}
\theoremstyle{definition}
\newtheorem{definition}[theorem]{Definition}
\theoremstyle{remark}
\newtheorem{remark}[theorem]{Remark}
\theoremstyle{plain}
\newtheorem{conjecture}[theorem]{Conjecture}
\theoremstyle{definition}
\newtheorem{axiom}[theorem]{Axiom}

\title{Geometric Origin of Nucleon Binding\\and a Programme for Beta Decay in HQIV\\{\large\scshape (working copy)}}
\author{Steven Ettinger Jr}
\date{\today}

\begin{document}
\maketitle

\begin{abstract}
We develop a self-contained geometric account of nucleon binding within the Horizon-Quantized Information Vacuum (HQIV) framework, built on a $3$D discrete null lattice with an octonionic eightfold mode count.
Cumulative simplex counting yields a shell-wise curvature-imprint ratio identically equal to $\alpha=3/5$; the informational monogamy complement is then $\gamma=1-\alpha=2/5$.
We record the synchronous-comoving lapse $N=1+\Phi+\phi t$ (natural units), the identification of the horizon auxiliary $\phi$ across the emergent gauge sector and an effective Friedmann coupling $G_{\mathrm{eff}}(\phi)=\phi^{\alpha}$, and a concrete \emph{divisor} prescription that renormalizes nucleon meta-harmonic masses on the lock-in shell when $N>1$.
Elementary inequalities show the lapse-corrected readout is strictly below the uncorrected mass in that regime.
Beta decay is proposed as a curvature-balance transition coupled to \emph{skew} time-phase alignment on the octonionic carrier; explicit lifetimes and a full nuclear potential are not derived here.
We close with a short comparison to Standard Model parameter counts and with an appendix that maps each major claim to an independent machine-checked certificate in a public repository---an audit layer that does not substitute for the in-text arguments.
\end{abstract}

\section{Discrete null lattice and the imprint identities}

\subsection{Mode counts}

Fix a non-negative shell index $m$.
The positive quadrant slice $\{(x,y,z)\in\mathbb{Z}_{\ge 0}^3 : x+y+z=m\}$ has $\binom{m+2}{2}$ integer points (stars-and-bars).
Following HQIV's octonionic carrier bookkeeping, attach a factor of eight to obtain the \emph{available} horizon modes at shell $m$,
\begin{equation}\label{eq:Am}
  A(m) \;=\; 8\binom{m+2}{2} \;=\; 4(m+1)(m+2)\,.
\end{equation}
Write $N_{\mathrm{simp}}(m):=(m+1)(m+2)$ for the \emph{stars-and-bars numerator} attached to shell $m$ in the HQIV convention (equal to twice a binomial coefficient, $2\binom{m+2}{2}$, if one insists on the classical $\binom{m+2}{2}$ point count).
The cumulative count up to $n$ is
\begin{equation}\label{eq:Cn}
  C(n) \;:=\; \sum_{m=0}^{n} N_{\mathrm{simp}}(m)\,.
\end{equation}

\begin{lemma}[Cubic cumulant / hockey-stick]\label{lem:hockey}
For every $n\in\mathbb{N}$,
\[
  3\,C(n) = (n+1)(n+2)(n+3)\,.
\]
Equivalently $C(n)=\dfrac{(n+1)(n+2)(n+3)}{3}$.
\end{lemma}

\begin{proof}
Induction on $n$.
For $n=0$, $C(0)=N_{\mathrm{simp}}(0)=2$ and $(0+1)(0+2)(0+3)=6=3C(0)$.
Assume $3C(n)=(n+1)(n+2)(n+3)$.
By definition $C(n+1)=C(n)+N_{\mathrm{simp}}(n+1)$ with $N_{\mathrm{simp}}(n+1)=(n+2)(n+3)$.
Thus
\begin{align*}
  3C(n+1) &= 3C(n)+3(n+2)(n+3) \\
  &= (n+1)(n+2)(n+3)+3(n+2)(n+3) \\
  &= (n+2)(n+3)\bigl((n+1)+3\bigr) \\
  &= (n+2)(n+3)(n+4)\,,
\end{align*}
which is the desired identity with $n$ replaced by $n+1$.
\end{proof}

\begin{definition}[Curvature-imprint ratio]\label{def:ratio}
For $n\in\mathbb{N}$ define
\[
  r(n) \;:=\; \frac{(n+1)(n+2)(n+3)}{5\,C(n)}\,.
\]
\end{definition}

\begin{theorem}[Forced imprint exponent]\label{thm:alpha}
For every $n\in\mathbb{N}$, $r(n)=3/5$.
Define $\alpha:=3/5$ and $\gamma:=1-\alpha=2/5$.
\end{theorem}

\begin{proof}
By Lemma~\ref{lem:hockey}, $(n+1)(n+2)(n+3)=3C(n)$, hence $r(n)=3C(n)/(5C(n))=3/5$ because $C(n)>0$ for all $n$.
\end{proof}

\begin{remark}
The rational $\alpha=3/5$ is therefore not a tunable exponent within this discrete normalization: it is the unique value of the ratio~\eqref{def:ratio} once~\eqref{eq:Cn} is adopted.
The companion $\gamma=2/5$ is the unique affine complement on the unit split $\alpha+\gamma=1$.
\end{remark}

\subsection{Temperature ladder and auxiliary field}

\begin{definition}[Shell temperature and auxiliary $\phi$]\label{def:phi}
Set $T(m):=1/(m+1)$ in natural units and $S(m):=(m+1)(m+2)$.
The HQIV homogeneous auxiliary field at shell $m$ is
\[
  \phi(m) \;:=\; \frac{2}{T(m)} \;=\; 2(m+1)\,.
\]
\end{definition}

\subsection{Electroweak scale witnesses (algebraic)}

The geometric vacuum expectation used in the HQIV electroweak package is assembled from the lock-in shell index $m_{\mathrm{lockin}}$, the ladder $T$, the surface factor $S$, and the monogamy complement $(1+\gamma)$; we do not reproduce the full triality-normalized coupling algebra in this note.
Specializing the closed rational outputs recorded in the development gives the numerical identities
\begin{equation}\label{eq:MW}
  M_W^{\mathrm{derived}}=\frac{392}{5}\,,\qquad
  m_H^{\mathrm{derived}}=\frac{588}{5}
\end{equation}
in the same mass units used for those lemmas, together with a companion neutral-boson relation tying $M_Z^{\mathrm{derived}}$ to $M_W^{\mathrm{derived}}$ and $\gamma$.

\section{Three naturals, one real, and skew time-phase alignment}

\subsection{Synthesis theorem (quoted package)}

Under the \emph{exact} quadratic shell law $A(m)=4(m+1)(m+2)$, transverse null-mode counting is modeled by three independent nonnegative integers $(x,y,z)\in\mathbb{N}_0^3$ with $x+y+z=m$ (the ``three naturals''), while causal height along maximal chains supplies an ordered timelike axis modeled by a single copy of $\mathbb{R}$ (the ``one real'').
The Fano-plane / octonion internal structure is then tied to the projective geometry of $\mathbb{F}_2^3$, and the imprint pair $(\alpha,\gamma)=(3/5,2/5)$ is the $d=3$ row of the standard HQIV unit-split family.
A fully self-contained statement and proof sketch of this \textbf{Three Natural + one Real} package appears in the companion HQIV manuscript on $3$D causal growth and octonionic gauge (Theorem ``Three Natural + one Real under the HQIV shell hypothesis'' there); the present note uses only the counting input~\eqref{eq:Am} and Theorem~\ref{thm:alpha}.

\subsection{Jobs of $\phi$ and $N$, and a fourth skew job}

Already in the HQVM chart one assigns three roles to the auxiliary field $\phi$ and the lapse $N=1+\Phi+\phi t$ (with $c=1$): \emph{(i)}~temperature / shell ladder ($T(m)=1/(m+1)$, $\phi(m)=2(m+1)$); \emph{(ii)}~hypercharge-like $\mathrm{U}(1)$ phase data via the octonionic scalar slot and the distinguished phase-lift generator $\Delta\in\mathfrak{so}(8)$ acting on $\mathbb{R}^8$; \emph{(iii)}~superposition bookkeeping on the $\mathbb{R}^8$ carrier.
We now add a fourth, explicitly \emph{geometric} role: \textbf{skew alignment} of the internal phase with the real timelike axis selected by the lapse.

The generator $\Delta$ is antisymmetric (skew-adjoint) in the $(e_1,e_7)$ plane in the standard octonion matrix model; octonion multiplication supplies non-associativity and a natural associator $3$-form that can be read as an algebraic ``skew'' coupling between the three transverse null directions and the ordered time axis.
The proposal is that nuclear binding and decay channels are selected not only by energy and charge, but by whether a given resonance product can \emph{close} curvature and \emph{realign} this skew factor inside a deep potential well.

\begin{axiom}[Skew time-phase alignment]\label{ax:skew}
Fix a skew-adjoint endomorphism $\Delta\in\mathfrak{so}(8)$ representing the phase-lift direction on $\mathbb{R}^8\cong\mathbb{O}$.
Let $s(t)>0$ denote the positive scalar factor assembled from the horizon auxiliary and lapse readout along the comoving slice (reducing to $\phi\,t$ in the homogeneous sector up to normalization choices), and let $\theta:\mathbb{R}\to\mathbb{R}$ be an angular coordinate.
The internal time-phase factor is modeled by
\begin{equation}\label{eq:skew-phase}
  \Psi(t) \;:=\; s(t)\,\exp\!\bigl(\theta(t)\,\Delta\bigr)\,\in\,\mathbb{R}_{>0}\times \mathrm{SO}(8)\,,
\end{equation}
with $\dot{\theta}$ proportional to the cumulative lapse increment along the chosen comoving observer in normalized units.
\end{axiom}

\begin{remark}[Status]
Equation~\eqref{eq:skew-phase} is a \textbf{modeling axiom}: it packages the idea that the lapse-selected ``now'' carries not only scalar dragging but an $\mathrm{SO}(8)$ phase generated by the same $\Delta$ that already appears in the $G_2\cup\{\Delta\}\Rightarrow\mathfrak{so}(8)$ closure story.
What is \emph{proved} in the discrete layer are the shell identities and Theorem~\ref{thm:alpha}; what is \emph{proved} in the matrix layer is skew-adjointness of the concrete $\Delta$ matrix used in the repository.
Wiring~\eqref{eq:skew-phase} into a variational principle for nuclear wells or $\beta$ widths is left open.
\end{remark}

\section{Lapse, gravity coupling, and nucleon divisor readouts}

\subsection{Synchronous lapse and shared $\phi$}

\begin{definition}[ADM lapse in the HQVM chart]\label{def:lapse}
Let $\Phi$ be the Newtonian potential, $\phi$ the horizon auxiliary, and $t$ coordinate time (natural units, $c=1$).
Define the lapse
\begin{equation}\label{eq:lapse}
  N(\Phi,\phi,t) \;:=\; 1+\Phi+\phi\,t\,.
\end{equation}
\end{definition}

\begin{definition}[Effective gravitational coupling]\label{def:Geff}
For $\phi\ge 0$ set $G_{\mathrm{eff}}(\phi):=\phi^{\alpha}$ with $\alpha=3/5$ from Theorem~\ref{thm:alpha}.
\end{definition}

\begin{remark}[Compatibility package]
The HQIV construction posits that the \emph{same} real scalar $\phi$ enters the lapse time-angle term $\phi t$, an emergent O-type Maxwell sector, and the Friedmann-side coupling $G_{\mathrm{eff}}$.
The precise variational derivation of Einstein equations from matter is \emph{not} claimed here; what is used below is the algebraic synchronization of symbols across sectors.
\end{remark}

\subsection{Lock-in shell and binding readout}

\begin{definition}[Lock-in index]
$m_{\mathrm{lockin}}\in\mathbb{N}$ denotes the discrete shell at which QCD-scale structure is aligned with horizon lock-in in the HQIV ladder.
It is treated as a fixed structural index in the same sense as a renormalization scale in ordinary field theory: downstream masses are \emph{evaluated} there, not solved for from~\eqref{eq:Am} alone.
\end{definition}

Let $B_{\star}>0$ denote the shared binding energy (in mass units) extracted from the composite-trace functional at $m_{\mathrm{lockin}}$.
Let $M_u,M_d>0$ be dressed constituent masses for the light quarks in the same units, assembled into proton and neutron \emph{constituent} totals
\[
  M_p^{\mathrm{const}} := 2M_u+M_d\,,\qquad M_n^{\mathrm{const}} := M_u+2M_d\,.
\]
Define the \emph{raw} meta-harmonic nucleon masses by subtraction of the shared binding,
\begin{equation}\label{eq:rawmass}
  M_p^{\mathrm{raw}} := M_p^{\mathrm{const}}-B_{\star}\,,\qquad
  M_n^{\mathrm{raw}} := M_n^{\mathrm{const}}-B_{\star}\,.
\end{equation}

On the lock-in shell set $\phi_{\star}:=\phi(m_{\mathrm{lockin}})$ and define the restricted lapse $N_{\star}(\Phi,t):=N(\Phi,\phi_{\star},t)$.

\begin{proposition}[Lapse divisor correction]\label{prop:divisor}
Define corrected masses by
\[
  M_p^{\mathrm{corr}}(\Phi,t):=\frac{M_p^{\mathrm{raw}}}{N_{\star}(\Phi,t)}\,,\qquad
  M_n^{\mathrm{corr}}(\Phi,t):=\frac{M_n^{\mathrm{raw}}}{N_{\star}(\Phi,t)}\,.
\]
If $N_{\star}(\Phi,t)>0$, the prescriptions are well-defined.
\end{proposition}

\begin{proof}
Immediate from~\eqref{eq:lapse} and division on $\mathbb{R}$.
\end{proof}

\begin{lemma}[Positivity of $\phi_{\star}$]\label{lem:phipos}
For every shell index $m$, $\phi(m)>0$.
In particular $\phi_{\star}>0$.
\end{lemma}

\begin{proof}
By Definition~\ref{def:phi}, $\phi(m)=2(m+1)>0$.
\end{proof}

\begin{proposition}[Weak-field forward-time regime]\label{prop:Npos}
If $1+\Phi>0$, $\phi=\phi_{\star}$, and $t\ge 0$, then $N(\Phi,\phi,t)>0$.
If in addition $\Phi\ge 0$ and $t>0$, then $N(\Phi,\phi,t)>1$.
\end{proposition}

\begin{proof}
From~\eqref{eq:lapse}, $N=1+\Phi+\phi t\ge 1+\Phi>0$ when $t\ge 0$ and $\phi>0$ (Lemma~\ref{lem:phipos}).
If $\Phi\ge 0$ and $t>0$, then $\phi t>0$ and hence $N>1+\Phi\ge 1$, in fact $N>1$ because $\phi t>0$.
\end{proof}

\begin{corollary}[Mass readout strictly decreases when $N>1$]\label{cor:mono}
If $M_p^{\mathrm{raw}}>0$ and $N_{\star}(\Phi,t)>1$, then $M_p^{\mathrm{corr}}(\Phi,t)<M_p^{\mathrm{raw}}$.
The same statement holds with $M_n$ in place of $M_p$.
\end{corollary}

\begin{proof}
For $N>1$ and $M>0$, $M/N<M$ is equivalent to $M<MN$, i.e.\ to $1<N$ after cancelling $M$.
\end{proof}

\subsection{Isospin splitting (model input vs.\ data)}

The HQIV nucleon package separates a fixed hypercharge sign contribution (normalized to a $1\,\mathrm{MeV}$ bookkeeping gap in the formal witness) from a residual mass difference driven by the dressed light-quark average and detuning dress terms.
The observed neutron--proton mass difference $\approx 1.293\,\mathrm{MeV}$ is quoted only as a phenomenological benchmark: matching it quantitatively is not asserted in this working copy.

\begin{remark}[PDG gap]
Current meta-harmonic nucleon outputs in the computational pipeline sit several percent above PDG masses; closing that gap is stated as an open strong-sector and dressing problem, not as a defect of Theorem~\ref{thm:alpha}.
\end{remark}

\section{A programme for beta decay}

\begin{conjecture}[Curvature and skew: $\beta$ transitions]\label{conj:beta}
Outside compact nuclear environments, a free neutron occupies a branch of the curvature ledger with sub-lock-in normalized curvature ($\Omega_k$-type ratios below unity relative to the calibration shell) and with \textbf{skew misalignment} of the internal phase~\eqref{eq:skew-phase} relative to the real timelike axis.
Embedding in a nuclear potential well together with lapse reinforcement is conjectured to supply both sufficient \emph{depth} and a triality-compatible \emph{skew realignment}, closing curvature only for resonance products whose Spin$(8)$ triality channel matches the local skew configuration (charge and magnetic moment are partial selectors; skew completes the orientation).
Conversely, proton-rich configurations whose wells \emph{overshoot} in skew-depth are conjectured to reopen the imbalance and relax via $\beta^{+}$ or electron-capture channels---the inverse geometric channel to neutron stabilization.
\end{conjecture}

\begin{remark}[Scope]
Conjecture~\ref{conj:beta} is deliberately schematic: no Fermi theory, no $W$ propagator, and no partial width is derived.
The role of the weak sector in HQIV is instead represented by structural reductions of an octonionic weak--Higgs scaffold onto an abelian Maxwell-type backbone, together with triality bookkeeping that enforces a neutral CP-tilt average and a constructive spin--statistics package compatible with horizon-limited continuum scaffolds.
Because $\Delta$ acts on the full $\mathbb{R}^8$ carrier, the neutron's nonzero magnetic moment participates in skew/triality selection even though the neutron is electrically neutral at the $\mathrm{U}(1)$ hypercharge level.
\end{remark}

\section{Large-scale signatures}

\begin{remark}[Low multipoles]
If the lapse-selected time slicing couples coherently to the longest-wavelength modes of the background, one expects modified low-$\ell$ statistics in polarization and temperature (sometimes colloquially summarized as ``axis-of-evil''-type alignments).
Any such prediction must eventually pass through a full Einstein--Boltzmann transport implementation; that layer is not part of the present note.

Separately, a purely algebraic \emph{theory hook} relates a birefringence angle defined from $(\alpha,m_{\mathrm{lockin}})$ to a redshift factor $(1+z)=(m_{\mathrm{lockin}}+1)^{\alpha}$ at a chosen normalization $\kappa_{\beta}=1$.
This identity is a consistency display, not a claim that the cosmic birefringence angle has been measured to equal that value.
\end{remark}

\section{Comparison to the Standard Model}

The electroweak sector of the Standard Model introduces an independent non-abelian weak gauge structure and, in textbook counts, on the order of twenty-six freely fitted parameters.
The HQIV presentation here reduces a portion of that freedom to discrete counting, a single imprint exponent forced by Lemma~\ref{lem:hockey}, and explicit structural pins (lock-in index, quark-sector anchors) whose remaining arbitrariness is acknowledged rather than hidden.

\section*{Acknowledgements}
The author acknowledges AI-assisted drafting and code navigation; mathematical claims in the main text are intended to be verifiable by hand from the definitions given.

\medskip
\noindent\textit{Further reading (HQIV programme).} \cite{ettinger2026hqiv,brodie2026}

\appendix

\section{Formal audit layer (optional)}

The repository \repo{} contains a Lean~4 formalization used as an \emph{independent check} of the algebraic claims above: Theorem~\ref{thm:alpha} is the statement \texttt{latticeAlphaRatio\_eq\_alpha} specialized to the definition of $\alpha$; Lemma~\ref{lem:phipos} is \texttt{phi\_of\_shell\_pos}; Proposition~\ref{prop:Npos} re-packages \texttt{HQVM\_lapse\_pos} and \texttt{HQVM\_lapse\_gt\_one}; Corollary~\ref{cor:mono} is \texttt{protonMassFromMetaHarmonics\_lapseCorrected\_lt\_raw\_of\_lapse\_gt\_one}.
The divisor law Proposition~\ref{prop:divisor} is \texttt{protonMassFromMetaHarmonics\_lapseCorrected\_eq\_raw\_div\_lapse}.
Skew-adjointness of the concrete phase-lift matrix for $\Delta$ is \texttt{phaseLiftDelta\_antisymm} in \texttt{Hqiv/Algebra/PhaseLiftDelta.lean} (re-exported through the SO$(8)$ / QFT feed modules).
For patch-local digital evolution, \texttt{Hqiv/QuantumComputing/HamiltonianToGateMapping.lean} formalizes sequential gate schedules with per-tick $\eta_\phi$--CKW certificates (\texttt{SequentialHQIVEvolution}), dense execution \texttt{runSequentialEvolution} (aligned with \texttt{DiscreteSchrodinger}), and an OSHoracle sparse trajectory \texttt{runSparseSequentialOSH}; these layers audit one-step-at-a-time Hamiltonian \emph{proxies} on the discrete cone without changing the shell-counting inputs used above.
Axiom~\ref{ax:skew} is \textbf{not} a Lean theorem: it records the modeling target coupling $\Delta$ to the lapse-selected slice.
No result in this appendix substitutes for the proofs in Sections~1--2: it is an audit trail for readers who wish to replay the certificates mechanically.

\section*{References}
\bibliographystyle{plain}
\bibliography{../references}

\end{document}
